\documentclass[../../main.tex]{subfiles}
\begin{document}

\begin{center}
  {\fontsize{48pt}{58pt}\selectfont \textbf{1993}\normalsize\textbf{年度}}
\end{center}

\begin{table}[htb]
  \centering
  \caption{出題テーマと難易度}
  \begin{tabular}{|c|p{9cm}|c|}
    \hline
    \textbf{問題番号} & \textbf{テーマ} & \textbf{難易度} \\
    \hline
    第1問 & 四面体の体積の極限 & \\
    \hline
    第2問 & 漸化式で定まる数列の整除性 & \\
    \hline
    第3問 & 障害物を避ける最短経路と軌跡 & \\
    \hline
    第4問 & 積分の最小値と最適な係数 & \\
    \hline
    第5問 & 書き写しの誤りと確率の極限 & \\
    \hline
    第6問 & 媒介変数で表された曲線と速度 & \\
    \hline
  \end{tabular}
\end{table}

\begin{table}[htb]
  \centering
  \caption{解答}
  \renewcommand{\arraystretch}{1.4}
  \begin{tabular}{|c|c|p{9cm}|}
    \hline
    \textbf{問題番号} & \textbf{小問} & \textbf{答え} \\
    \hline
    第1問 & -- & 極限値は$8$ \\
    \hline
    第2問 & (1) & $a_n$が偶数となるのは$n$が$3$の倍数のとき，かつそのときに限る \\
    \cline{2-3}
    & (2) & $a_n$が$10$の倍数となるのは$n$が$12$の倍数のとき，かつそのときに限る \\
    \hline
    第3問 & -- & 軌跡は解答の図を参照 \\
    \hline
    第4問 & -- & $n$が奇数なら$(p,q)=\left(-\dfrac{3}{n+2},0\right)$，$I_{\min}=\dfrac{(n-1)^2}{(2n+1)(n+2)^2}$．$n$が偶数なら$(p,q)=\left(0,-\dfrac{1}{n+1}\right)$，$I_{\min}=\dfrac{n^2}{(2n+1)(n+1)^2}$ \\
    \hline
    第5問 & -- & $\displaystyle\lim_{n\to\infty}C(n)=\dfrac{p^3q^2}{(p+q)^5}$ \\
    \hline
    第6問 & (1) & 最大値は$\sqrt{13}$，最小値は$\sqrt{\dfrac{8}{27}}$ \\
    \cline{2-3}
    & (2) & 自己交点は$(-1,0)$のみ．速度ベクトルは$(-2,-2)$，$(0,2)$，$(2,-2)$（図示は解答参照） \\
    \hline
  \end{tabular}
\end{table}

\paragraph{試験形態}
6問．試験時間は要確認．

\end{document}
